% ============================================================
% MERGED SILIQUN JOURNAL PAPER — SKELETON + SECTION MAPPING
% Title: SiliQun: A Software-Defined Quantum Firmware Framework
%        for Silicon Spin Qubits with Deep Reinforcement Learning
% Target: npj Quantum Information (primary) / QST (secondary) /
%         IEEE TQE (fallback)
%
% SOURCE MAPPING KEY:
%   [P1]  = QUASAR_Paper1_Full_Manuscript.md (SiliQun simulator paper)
%   [P6]  = paper6-siliqun-framework/*.tex   (SiliQun framework/PGIRS paper)
%   [NEW] = New content required for the merged manuscript
% ============================================================

\documentclass[twocolumn]{article}
\usepackage{amsmath,amssymb,physics,graphicx,booktabs,hyperref,cleveref}
\usepackage[margin=1in]{geometry}

% ── TITLE ────────────────────────────────────────────────────────────────────
% [NEW] Merged title that foregrounds both the framework AND the result.
% Rationale: npj QI and QST reward papers that introduce a reusable artifact
% AND demonstrate its scientific utility in a single manuscript.
\title{SiliQun: A Software-Defined Quantum Firmware Framework for Silicon
       Spin Qubits with Deep Reinforcement Learning Validation}

\author{Raad A. Al-Shehri$^{1}$ \\
        $^{1}$Faculty of Computing and Information Technology,
        King Abdulaziz University, Jeddah, Saudi Arabia \\
        \texttt{ralshehri0468@stu.kau.edu.sa}}

\date{}

\begin{document}
\maketitle

% ── ABSTRACT ─────────────────────────────────────────────────────────────────
% SOURCE: Merge of P6 abstract (framework framing) + P1 abstract (result claim)
% STRATEGY: Lead with the SDQF concept (P6), then land the F≥0.99 result (P1).
% LENGTH TARGET: 250 words (npj QI limit).
%
% [P6] → paragraphs 1–3: firmware gap, SiliQun architecture, PGIRS methodology,
%         Gymnasium interface, Lindblad physics, analytical validation benchmarks
% [P1] → paragraph 4: QUASAR v9.51 result — F≥0.99 on 4q GHZ under SiMOS noise,
%         BRFD reward discovery, DEHB hyperparameter governance, expert warm-start
% [P6] → paragraph 5: cross-algorithm benchmark (7 DRL families, F=0.847–0.993),
%         cross-backend validation via QuantumExecutor, MIT licence
% [NEW] → closing sentence: positions SiliQun as the first open-source SDQF layer
%         for SiMOS — analogous to BIOS/UEFI in classical computing
\begin{abstract}
% TODO: Write merged abstract (~250 words) following the source mapping above.
% Draft hook: "Between physical quantum hardware and high-level quantum algorithms
% lies a critical but under-theorised abstraction layer — the quantum firmware gap.
% We present SiliQun, an open-source framework that closes this gap for silicon
% metal-oxide-semiconductor (SiMOS) spin qubits by implementing the Software-Defined
% Quantum Firmware (SDQF) architecture..."
\end{abstract}

% ── §1 INTRODUCTION ──────────────────────────────────────────────────────────
% SOURCE MAPPING:
%   [P6 §1, ¶1–3]  → The firmware gap concept; SDQF as BIOS/UEFI analogy;
%                     control/data plane separation (SiliQun=data, QUASAR=control)
%   [P1 §1, ¶1–4]  → The fault-tolerance threshold motivation; why SiMOS is the
%                     target; the gap in existing RL environments (abstract noise)
%   [NEW ¶]        → Explicit statement of the dual contribution: (1) framework
%                     infrastructure, (2) empirical validation crossing F≥0.99
%   [P6 §1, ¶last] → Contributions list (5 items from P6, augmented with P1 result)
%   [P1 §1, ¶last] → Paper organisation paragraph (adapt for merged structure)
%
% WRITING NOTE: The introduction must not read as two papers glued together.
% The unifying narrative is: "We built the firmware layer (SiliQun/PGIRS),
% then used it to demonstrate the first fault-tolerance-threshold result on
% a physically calibrated SiMOS simulator." The framework enables the result;
% the result validates the framework.
\section{Introduction}
% TODO: ~800 words. Follow source mapping above.

% ── §2 RELATED WORK ──────────────────────────────────────────────────────────
% SOURCE MAPPING:
%   [P6 §2, all subsections] → USE AS-IS from sec2_related_work.tex
%     §2.1 RL Environments for Quantum Circuit Synthesis
%     §2.2 Pulse-Level Quantum Control with RL
%     §2.3 Quantum Noise Simulation Frameworks
%     §2.4 The Closest Competitor: Orbell (2024) Oxford DPhil
%     §2.5 Fault-Tolerant Logical Qubit Encoding with RL
%     §2.6 Summary: SiliQun's Unique Contribution
%   [P1 §2]  → MERGE selectively: P1's related work covers RL for gate-based
%              circuit synthesis (Moro, He, Kuo, Bi) — these belong in §2.1
%              of the merged paper. They are NOT currently in P6 §2.1.
%              ACTION: Extract P1 §2 ¶1–4 and insert into merged §2.1 as the
%              "gate-level DRL" sub-paragraph before the environment papers.
%
% NET RESULT: Merged §2 = P6 §2 + P1 §2 gate-synthesis paragraph inserted into §2.1
\section{Related Work}
\input{sec2_related_work}  % P6 base — already written

% ── §3 THE SILIQUN FRAMEWORK ─────────────────────────────────────────────────
% SOURCE MAPPING:
%   [P6 §3 — NOT YET WRITTEN] → This is the main framework section.
%   Content to write (new for merged paper, expanding P6 §3 skeleton):
%
%   §3.1 The SDQF Architecture and Firmware Gap
%     [P6 §1 intro material] → Expand the firmware gap concept into a full subsection.
%     [NEW] → Formal definition of SDQF: control plane (QUASAR/DRL agents) +
%             data plane (SiliQun Lindblad physics engine) + interface layer (Gymnasium)
%
%   §3.2 Physics Engine: Lindblad Master Equation Solver
%     [P1 §3.2 SiMOS environment] → The SiMOS simulation environment description.
%     [P6 abstract/intro physics paragraphs] → T1 amplitude damping, T2 dephasing,
%       depolarising noise, SiMOS-calibrated parameters (T1=1ms, T2=0.5ms, p2=1%)
%     [NEW] → Formal Lindblad equation with all three collapse operators;
%             parameter table (T1, T2, gate fidelity, qubit frequency)
%
%   §3.3 Gymnasium Interface and Environment Design
%     [P6 §3 — to be written] → Action space (discrete gate set: H, CNOT, Rz, CZ),
%       observation space (density matrix flattened), reward structure,
%       episode termination conditions, reset protocol
%     [P1 §3.1 Problem Formulation] → MDP formulation (S, A, R, γ, T)
%
%   §3.4 The PGIRS Methodology
%     [P6 §3 PGIRS — to be written] → Five-phase engineering discipline:
%       Phase 1: Noise-first design
%       Phase 2: Interface abstraction
%       Phase 3: Reward engineering (BRFD)
%       Phase 4: Hierarchical decomposition
%       Phase 5: Empirical validation
%
%   §3.5 Analytical Validation Benchmarks
%     [P6 abstract validation results] → Three benchmarks:
%       (a) Rabi oscillations: max deviation 0.00 (machine precision)
%       (b) T1 amplitude decay: max deviation 2.78×10⁻¹⁶ (float64 epsilon)
%       (c) Bell state fidelity vs coherence time: max deviation 0.0024
%     [NEW] → Figures: Rabi oscillation plot, T1 decay curve, Bell fidelity curve
\section{The SiliQun Framework}

\subsection{The SDQF Architecture and the Firmware Gap}
% TODO: ~400 words. Sources: [P6 §1] + [NEW formal SDQF definition]

\subsection{Physics Engine: Lindblad Master Equation Solver}
% TODO: ~500 words. Sources: [P1 §3.2] + [P6 physics paragraphs] + [NEW Lindblad eq]

\subsection{Gymnasium Interface and MDP Formulation}
% TODO: ~400 words. Sources: [P1 §3.1] + [P6 §3 interface — to write]

\subsection{The PGIRS Methodology}
% TODO: ~500 words. Sources: [P6 §3 PGIRS — to write]

\subsection{Analytical Validation Benchmarks}
% TODO: ~400 words + 3 figures. Sources: [P6 abstract validation] + [NEW figures]

% ── §4 QUASAR v9.51: DRL AGENT ARCHITECTURE ──────────────────────────────────
% SOURCE MAPPING:
%   [P1 §3.3] → QUASAR v9.51 Architecture: 7 specialised training components,
%               event-driven registry, SAC algorithm choice
%   [P1 §3.4] → Curriculum Learning and Expert Injection (warm-start)
%   [P1 §3.3] → BRFD: Bayesian Reward Function Discovery
%   [P1 §3.3] → DEHB: Differential Evolution Hyperband governance
%   [NEW]     → Brief paragraph positioning QUASAR as the "control plane" agent
%               that runs on top of the SiliQun "data plane" — making the
%               SDQF architecture concrete
%
% WRITING NOTE: This section is essentially P1 §3 reframed as "here is the DRL
% agent we used to validate SiliQun" rather than "here is our main contribution."
% The framing shift is subtle but important for the merged paper's narrative.
\section{QUASAR v9.51: The DRL Control Plane}

\subsection{Agent Architecture and Training Components}
% TODO: ~600 words. Source: [P1 §3.3]

\subsection{Autonomous Reward Discovery via BRFD}
% TODO: ~300 words. Source: [P1 §3.3 BRFD subsection]

\subsection{Hyperparameter Governance via DEHB}
% TODO: ~300 words. Source: [P1 §3.3 DEHB subsection]

\subsection{Curriculum Learning and Expert Warm-Start}
% TODO: ~300 words. Source: [P1 §3.4]

% ── §5 EXPERIMENTAL VALIDATION ───────────────────────────────────────────────
% SOURCE MAPPING — this is the most complex merge section:
%
%   §5.1 Experimental Setup
%     [P1 §4.1] → Hardware (NVIDIA RTX 2070), software stack (PyTorch, SB3),
%                 training budget (500k steps), evaluation protocol (100 episodes)
%     [NEW]     → Add: SiliQun version, Gymnasium version, Python version,
%                 reproducibility statement (seed, GitHub link)
%
%   §5.2 Primary Result: Fault-Tolerance Threshold on 4q GHZ
%     [P1 §4.2] → Fidelity progression and convergence (F=0.9923 best, F≥0.99 threshold)
%     [P1 §4.3] → Comparison with baseline pipeline (ablation of BRFD, DEHB, warm-start)
%     [P1 §4.4] → Component contribution analysis
%     [P1 §4.5] → Statistical significance (5 seeds, mean±std)
%     [P1 §4.6] → Fault-tolerance implications (threshold theorem context)
%
%   §5.3 Cross-Algorithm Benchmark: Seven DRL Families
%     [P6 abstract cross-algorithm results] → F=0.847 (DQN, T3) to F=0.993 (SAC-HMRL, T7)
%     [NEW]     → Full table: 7 algorithm families × fidelity × convergence steps
%                 Algorithms: DQN (T3), PPO (T4), SAC (T5), TD3 (T6), SAC-HMRL (T7),
%                 DDPG (T1), A2C (T2) — map to the 7-category taxonomy from P5 survey
%
%   §5.4 Cross-Backend Policy Transferability
%     [P6 §5.3 sec53_cross_backend_validation.tex] → USE AS-IS
%       QuantumExecutor VirtualProvider, 3 circuits × 4 backends × 4096 shots,
%       XEB + TVD metrics, transferability results
%
%   §5.5 Baseline Transferability Evaluation
%     [P1 §4.7 E5] → Transfer of trained policy to unseen noise configurations
\section{Experimental Validation}

\subsection{Experimental Setup}
% TODO: ~300 words. Sources: [P1 §4.1] + [NEW reproducibility statement]

\subsection{Primary Result: Crossing the Fault-Tolerance Threshold}
% TODO: ~700 words + 3 figures. Sources: [P1 §4.2–4.6]

\subsection{Cross-Algorithm Benchmark: Seven DRL Families}
% TODO: ~400 words + 1 table. Sources: [P6 abstract] + [NEW full table]

\subsection{Cross-Backend Policy Transferability}
% TODO: ~400 words + 1 table. Sources: [P6 §5.3 — already written]
\input{sec53_cross_backend_validation}  % P6 §5.3 — already written

\subsection{Noise Robustness and Baseline Transferability}
% TODO: ~300 words. Source: [P1 §4.7]

% ── §6 DISCUSSION ────────────────────────────────────────────────────────────
% SOURCE MAPPING:
%   [P1 §5.1] → Interpretation of the breakthrough result (F≥0.99 significance)
%   [P1 §5.2] → Comparison with related work (He et al., Bi et al., Orbell 2024)
%   [P6 §2.4 Orbell comparison] → Expand the Oxford DPhil comparison here
%   [NEW]     → Discussion of PGIRS as a reusable methodology beyond SiMOS:
%               applicability to trapped ions, superconducting qubits, NV centres
%   [NEW]     → SDQF as a software engineering paradigm: versioning, hot-patching,
%               reconfigurability analogies with classical firmware
%   [P1 §5.3] → Limitations: 4-qubit ceiling, SiMOS-only calibration,
%               simulation-to-hardware gap
\section{Discussion}

\subsection{Significance of the Fault-Tolerance Threshold Result}
% TODO: ~400 words. Source: [P1 §5.1]

\subsection{Positioning Against Related Work}
% TODO: ~500 words. Sources: [P1 §5.2] + [P6 §2.4 Orbell]

\subsection{Generalisability of the PGIRS Methodology}
% TODO: ~400 words. Source: [NEW]

\subsection{Limitations and Future Work}
% TODO: ~400 words. Sources: [P1 §5.3] + [NEW simulation-to-hardware gap discussion]

% ── §7 CONCLUSION ────────────────────────────────────────────────────────────
% SOURCE MAPPING:
%   [P1 §5.4] → Core conclusion paragraph (F≥0.99 result)
%   [P6 §1 last paragraph] → Framework conclusion (SDQF, open-source, PGIRS)
%   [NEW]     → Forward-looking sentence: ANDROMEDA (P3) as the next layer
%               (hierarchical modular RL for logical qubit construction on SiliQun)
\section{Conclusion}
% TODO: ~300 words. Sources: [P1 §5.4] + [P6 §1] + [NEW ANDROMEDA forward reference]

% ── REFERENCES ───────────────────────────────────────────────────────────────
% SOURCE MAPPING:
%   [P1 References]  → All P1 references (gate synthesis, SiMOS, BRFD, DEHB)
%   [P6 §2 bibitem blocks] → All P6 references (environments, Orbell, QuTiP, etc.)
%   [NEW]            → QuantumExecutor (arXiv 2507.07597v1), PGIRS methodology refs
%
% ACTION: Merge both reference lists, deduplicate, renumber sequentially.
% Estimated total: ~50–60 references (appropriate for npj QI / QST)
\bibliographystyle{unsrt}
\bibliography{merged_siliqun_refs}

\end{document}
